% titel: Blatt 6, A8: Molekülgeometrie nach VSEPR
% skript: Mortimer Kap. 9 (chemie: VSEPR-Modell, Molekülgeometrie)

Die Molekülgeometrie (Position der Atome) ergibt sich aus der
Elektronenpaar-Anordnung um das Zentralatom nach Abzug der freien Paare.
Freie Paare bestimmen die Anordnung, zählen aber für die Molekülgeometrie
nicht mit.

\begin{werkzeugkasten}[Geometrie bei 4--6 Elektronenpaaren]
\begin{center}
\begin{tabular}{@{}c l@{}}
  EP-Gebiete um Zentralatom & mögliche Molekülgeometrien \\
  \hline
  3  (1 frei) & gewinkelt (bent) \\
  4  (1 frei) & trigonal-pyramidal; (2 frei) gewinkelt/V-förmig \\
  5  (1 frei) & see-saw (quadratisch-/sägeförmig); (2 frei) T-förmig; \\
     & (3 frei) linear \\
  6  (1 frei) & quadratisch-pyramidal; (2 frei) quadratisch-planar \\
\end{tabular}
\end{center}
Zur Bestimmung: EP-Paare zählen (Bindungen + freie Paare), Anordnung aus
Abstossung ableiten, dann freie Paare ``weglassen'', um die Molekülgeometrie zu
erhalten.
\textit{Skript:} Mortimer Kap. 9 (VSEPR).
\end{werkzeugkasten}

\begin{achtung}[Darauf achten]
\begin{itemize}
  \item Unterschied: \textit{Elektronenpaar-Anordnung} (mit freien Paaren) vs.\
        \textit{Molekülgeometrie} (ohne). Beide klar benennen.
  \item Periode-3/4/5-Zentralatome (As, Br, I, Sn, Xe) haben oft $d$-Orbitale
        und damit 5--6 EP-Gebiete.
  \item Ladungen berücksichtigen: $\mathrm{XeF_3^+}$ hat wegen der Kationen-
        Ladung ein Valenzelektron weniger als neutrales Xe (eine $e^-$ weniger);
        das reduziert Zahl der freien Paare.
\end{itemize}
\end{achtung}
